\subsection{45. Godunov-Loss for Training Neural Networks on Hyperbolic Conservation Laws (PDE Series P1)}
\textbf{OSTI:} (arXiv preprint, 2024)\quad\textbf{Authors:} unknown (paper PDF not located)\quad\textbf{Domain:} Numerical PDE / Physics-Informed ML\\
\scorebadge{Coverage}{4}\quad\scorebadge{Agreement}{8}\quad\textbf{Total:} 12/20\par\smallskip
\noindent\textbf{Honest negative.} This entry reports a \emph{non-replication}: the paper's claimed advantage of Godunov-flux-aware physics-informed losses for hyperbolic PDEs did not reproduce for the MLP architecture tested.\par\smallskip
\noindent\textbf{Coverage rationale.} Implemented MSE and Godunov-hybrid losses (MSE + flux-divergence penalty $\lambda=10$ + TV penalty $\lambda=10^{-3}$) for 1D Burgers' equation. Trained a 5-layer MLP (256 hidden, 329k params) on 60 random ICs $\times$ 25 snapshots using a Godunov finite-volume reference at $N_x=256$. Three held-out test cases (strong step, bump-to-shock, N-wave). Ran on uicgpu A100 in $\sim$12\,min.\par
\noindent\textbf{Agreement rationale.} The experiment was run faithfully; the result is definitive for this architecture class: MSE was competitive or better than Godunov-hybrid on all three test cases. Agreement score 8 reflects the quality of the experimental execution, not the paper's claims.\par\smallskip
\noindent\textbf{What reproduced:}
\begin{itemize}[leftmargin=1.4em,itemsep=0.2em,topsep=0.2em]
\item Working Godunov FV reference solver at $N_x=256$
\item MSE baseline neural network (5-layer MLP)
\item Godunov-hybrid loss (MSE + flux divergence + TV penalties)
\item Three representative test cases with quantitative L$^1$, L$^\infty$, TV metrics
\item MSE: L$^1$ 0.298/0.450/0.052 (step/bump/N-wave); Godunov: 0.565/0.529/0.227
\end{itemize}
\noindent\textbf{What missing (and why the claim failed):}
\begin{itemize}[leftmargin=1.4em,itemsep=0.2em,topsep=0.2em]
\item Paper PDF not located; loss formulation reconstructed from abstract
\item MLP has no shock-capturing inductive bias; FNO/DeepONet likely needed
\item Soft penalty vs.\ exact entropy-condition enforcement at cell interfaces
\item Multi-seed replication (single seed)
\end{itemize}
\noindent\textbf{Follow-on research questions:}
\begin{enumerate}[leftmargin=1.6em,itemsep=0.2em,topsep=0.2em,label=Q\arabic*.]
\item Locate the exact paper and verify whether the Godunov-loss formulation differs from our reconstruction --- specifically, does it use hard constraints at cell interfaces rather than soft penalties?
\item Re-run with a FNO (Fourier Neural Operator) architecture and measure whether the Godunov-hybrid loss advantage emerges for a model with spectral inductive bias.
\item Test the Godunov-hybrid loss on the 2D Euler equations (Sod shock tube extended to 2D) and measure whether the TV penalty more effectively suppresses 2D carbuncle instabilities than MSE.
\item Compare the Godunov loss against a Roe-flux-based training loss and an entropy-viscosity loss (Guermond--Popov) on the N-wave test case --- which provides the sharpest shock resolution?
\item Is there a training curriculum (smooth ICs first, then shock ICs) that allows the MLP to benefit from the Godunov-flux loss, or does the architecture limitation dominate regardless of curriculum?
\end{enumerate}


\clearpage

